% --- Archon physics formalization source begin ---
% archon:physics
% archon:covers PhyXMiniProblems/problem_phyx_mini_0478.lean
% archon:source-report reports/phyx_mini/problem_phyx_mini_0478.source.json

\chapter{Physics problem phyx\_mini\_0478}
\label{ch:PhyXMiniProblems_problem_phyx_mini_0478}

\paragraph{Problem source.}
\#\# Physical scenario

Figure gives the lateral magnification \$m\$ of an object versus the object distance \$p\$ from a spherical mirror as the object is moved along the mirror’s central axis through a range of values for \$p\$. The horizontal scale is set by \$p\_s = 10.0 \textbackslash{}text\{cm\}\$.

\#\# Question

What is the magnification of the object when the object is 21 cm from the mirror?

\#\# Answer choices

- A: A: +0.16

- B: B: +0.22

- C: C: +0.28

- D: D: +0.32

\#\# Figure caption (auxiliary; use the image as primary evidence)

The image is a graph with a plot on a coordinate plane. 

- **Axes**: 
  - The horizontal axis is labeled \textbackslash{}( p \textbackslash{}) (cm), ranging from 0 to \textbackslash{}( p\_s \textbackslash{}).
  - The vertical axis is labeled \textbackslash{}( m \textbackslash{}), ranging from 0 to 1.

- **Curve**: A thick black curve starts at the top left, near \textbackslash{}( (0, 1) \textbackslash{}), and slopes downward to the right, closer to \textbackslash{}( (p\_s, 0) \textbackslash{}).

- **Grid**: The background is a grid to help visualize intervals.

- **Marks**: Minor tick marks are present on both axes, with the vertical axis specifically marked at approximately 0.5.

This graph seems to represent a decreasing relationship between \textbackslash{}( p \textbackslash{}) and \textbackslash{}( m \textbackslash{}).

\#\# Dataset metadata

Geometrical Optics; Physical Model Grounding Reasoning, Multi-Formula Reasoning

\paragraph{Recorded answer/context.}
D: D: +0.32

\paragraph{Figure/image path.}
/root/proposal\_for\_physic/science-mango/phyx\_mini\_run/phyx\_data/test\_image/478.png

\paragraph{Formalization target.}
create a compiling Lean file with sorry bodies at `PhyXMiniProblems/problem\_phyx\_mini\_0478.lean`. The Lean declarations must preserve the physical quantities, dimensions or dimensional roles, figure labels, governing-law hypotheses, and final relation expressed by this problem.
Use Mathlib/Physlib names found through LeanExplore where available. If a physics API is missing, introduce faithful local abstractions rather than scalar placeholder aliases.

\begin{theorem}[Physics formalization target]
\label{thm:physics:phyx_mini_0478:target}
\lean{PhyXMiniProblems.ProblemPhyXMini0478.problem_phyx_mini_0478}
\uses{def:physics:phyx-mini-0478:phyxminiproblems-problemphyxmini0478-mirrormagnificationexperiment, def:physics:phyx-mini-0478:phyxminiproblems-problemphyxmini0478-lateralmagnificationat, def:physics:phyx-mini-0478:phyxminiproblems-problemphyxmini0478-matchessphericalmirrorscenario, def:physics:phyx-mini-0478:phyxminiproblems-problemphyxmini0478-matchestextuallydescribedmagnificationgraph, def:physics:phyx-mini-0478:phyxminiproblems-problemphyxmini0478-hasintendedgraphcalibrationreadout, def:physics:phyx-mini-0478:phyxminiproblems-problemphyxmini0478-matchesquestionreadout, def:physics:phyx-mini-0478:phyxminiproblems-problemphyxmini0478-hasphysicalmirrorparameters, def:physics:phyx-mini-0478:phyxminiproblems-problemphyxmini0478-satisfiesparaxialsphericalmirrorlaws, def:physics:phyx-mini-0478:phyxminiproblems-problemphyxmini0478-answerchoice, def:physics:phyx-mini-0478:phyxminiproblems-problemphyxmini0478-matchesdisplayedmagnification}
The assigned autoformalize agent should translate this physics problem into Lean declarations in the covered file, with theorem and lemma proof bodies written as `by sorry`.
\end{theorem}
\begin{proof}
This is an autoformalization task, not a proof task. Produce statements that can later be proved by the physics prover without weakening the physical model.
\end{proof}

% --- Archon declaration topology begin ---
\section{Declaration topology}
\begin{definition}[Length Magnitude]
  \label{def:physics:phyx-mini-0478:phyxminiproblems-problemphyxmini0478-lengthmagnitude}
  \lean{PhyXMiniProblems.ProblemPhyXMini0478.LengthMagnitude}
  A nonnegative, unit-independent physical length magnitude
\end{definition}
\begin{proof}
  This is the declared model definition of Length Magnitude.
\end{proof}
\begin{definition}[Signed Optical Length]
  \label{def:physics:phyx-mini-0478:phyxminiproblems-problemphyxmini0478-signedopticallength}
  \lean{PhyXMiniProblems.ProblemPhyXMini0478.SignedOpticalLength}
  A signed, unit-independent optical length, used for f and q
\end{definition}
\begin{proof}
  This is the declared model definition of Signed Optical Length.
\end{proof}
\begin{definition}[unit Choices With Length Unit]
  \label{def:physics:phyx-mini-0478:phyxminiproblems-problemphyxmini0478-unitchoiceswithlengthunit}
  \lean{PhyXMiniProblems.ProblemPhyXMini0478.unitChoicesWithLengthUnit}
  SI base units with a selected unit replacing the SI length unit
\end{definition}
\begin{proof}
  This is the declared model definition of unit Choices With Length Unit.
\end{proof}
\begin{definition}[length Magnitude Readout]
  \label{def:physics:phyx-mini-0478:phyxminiproblems-problemphyxmini0478-lengthmagnitudereadout}
  \lean{PhyXMiniProblems.ProblemPhyXMini0478.lengthMagnitudeReadout}
  \uses{def:physics:phyx-mini-0478:phyxminiproblems-problemphyxmini0478-lengthmagnitude, def:physics:phyx-mini-0478:phyxminiproblems-problemphyxmini0478-unitchoiceswithlengthunit}
  Scalar readout of a nonnegative physical length in a selected unit
\end{definition}
\begin{proof}
  This is the declared model definition of length Magnitude Readout.
\end{proof}
\begin{definition}[signed Length Readout]
  \label{def:physics:phyx-mini-0478:phyxminiproblems-problemphyxmini0478-signedlengthreadout}
  \lean{PhyXMiniProblems.ProblemPhyXMini0478.signedLengthReadout}
  \uses{def:physics:phyx-mini-0478:phyxminiproblems-problemphyxmini0478-signedopticallength, def:physics:phyx-mini-0478:phyxminiproblems-problemphyxmini0478-unitchoiceswithlengthunit}
  Scalar readout of a signed optical length in a selected unit
\end{definition}
\begin{proof}
  This is the declared model definition of signed Length Readout.
\end{proof}
\begin{definition}[length In Centimeters]
  \label{def:physics:phyx-mini-0478:phyxminiproblems-problemphyxmini0478-lengthincentimeters}
  \lean{PhyXMiniProblems.ProblemPhyXMini0478.lengthInCentimeters}
  \uses{def:physics:phyx-mini-0478:phyxminiproblems-problemphyxmini0478-lengthmagnitude, def:physics:phyx-mini-0478:phyxminiproblems-problemphyxmini0478-lengthmagnitudereadout}
  Centimeter readout of a nonnegative physical length
\end{definition}
\begin{proof}
  This is the declared model definition of length In Centimeters.
\end{proof}
\begin{definition}[signed Length In Centimeters]
  \label{def:physics:phyx-mini-0478:phyxminiproblems-problemphyxmini0478-signedlengthincentimeters}
  \lean{PhyXMiniProblems.ProblemPhyXMini0478.signedLengthInCentimeters}
  \uses{def:physics:phyx-mini-0478:phyxminiproblems-problemphyxmini0478-signedopticallength, def:physics:phyx-mini-0478:phyxminiproblems-problemphyxmini0478-signedlengthreadout}
  Centimeter readout of a signed optical length
\end{definition}
\begin{proof}
  This is the declared model definition of signed Length In Centimeters.
\end{proof}
\begin{definition}[Mirror Geometry]
  \label{def:physics:phyx-mini-0478:phyxminiproblems-problemphyxmini0478-mirrorgeometry}
  \lean{PhyXMiniProblems.ProblemPhyXMini0478.MirrorGeometry}
  Geometry of the reflecting optical element
\end{definition}
\begin{proof}
  This is the declared model definition of Mirror Geometry.
\end{proof}
\begin{definition}[Object Motion Path]
  \label{def:physics:phyx-mini-0478:phyxminiproblems-problemphyxmini0478-objectmotionpath}
  \lean{PhyXMiniProblems.ProblemPhyXMini0478.ObjectMotionPath}
  Path along which the object is moved
\end{definition}
\begin{proof}
  This is the declared model definition of Object Motion Path.
\end{proof}
\begin{definition}[Optical Approximation]
  \label{def:physics:phyx-mini-0478:phyxminiproblems-problemphyxmini0478-opticalapproximation}
  \lean{PhyXMiniProblems.ProblemPhyXMini0478.OpticalApproximation}
  Optical model used for the mirror laws
\end{definition}
\begin{proof}
  This is the declared model definition of Optical Approximation.
\end{proof}
\begin{definition}[Graph Quantity]
  \label{def:physics:phyx-mini-0478:phyxminiproblems-problemphyxmini0478-graphquantity}
  \lean{PhyXMiniProblems.ProblemPhyXMini0478.GraphQuantity}
  Quantities named on the axes of the intended graph
\end{definition}
\begin{proof}
  This is the declared model definition of Graph Quantity.
\end{proof}
\begin{definition}[Magnification Distance Graph]
  \label{def:physics:phyx-mini-0478:phyxminiproblems-problemphyxmini0478-magnificationdistancegraph}
  \lean{PhyXMiniProblems.ProblemPhyXMini0478.MagnificationDistanceGraph}
  \uses{def:physics:phyx-mini-0478:phyxminiproblems-problemphyxmini0478-lengthmagnitude, def:physics:phyx-mini-0478:phyxminiproblems-problemphyxmini0478-graphquantity}
  Coordinate data for the intended graph. Horizontal coordinates are scalar centimeter readouts of physical object distances, whereas vertical coordinates are signed dimensionless magnifications
\end{definition}
\begin{proof}
  This is the declared model definition of Magnification Distance Graph.
\end{proof}
\begin{definition}[Mirror Magnification Experiment]
  \label{def:physics:phyx-mini-0478:phyxminiproblems-problemphyxmini0478-mirrormagnificationexperiment}
  \lean{PhyXMiniProblems.ProblemPhyXMini0478.MirrorMagnificationExperiment}
  \uses{def:physics:phyx-mini-0478:phyxminiproblems-problemphyxmini0478-lengthmagnitude, def:physics:phyx-mini-0478:phyxminiproblems-problemphyxmini0478-signedopticallength, def:physics:phyx-mini-0478:phyxminiproblems-problemphyxmini0478-mirrorgeometry, def:physics:phyx-mini-0478:phyxminiproblems-problemphyxmini0478-objectmotionpath, def:physics:phyx-mini-0478:phyxminiproblems-problemphyxmini0478-opticalapproximation, def:physics:phyx-mini-0478:phyxminiproblems-problemphyxmini0478-magnificationdistancegraph, def:physics:phyx-mini-0478:phyxminiproblems-problemphyxmini0478-sourcescene, def:physics:phyx-mini-0478:phyxminiproblems-problemphyxmini0478-evidencepolicy, def:physics:phyx-mini-0478:phyxminiproblems-problemphyxmini0478-suppliedrasteraudit}
  The physical mirror experiment. The focal length and the image-distance response are independent signed physical lengths; neither is defined from the requested answer
\end{definition}
\begin{proof}
  This is the declared model definition of Mirror Magnification Experiment.
\end{proof}
\begin{definition}[lateral Magnification At]
  \label{def:physics:phyx-mini-0478:phyxminiproblems-problemphyxmini0478-lateralmagnificationat}
  \lean{PhyXMiniProblems.ProblemPhyXMini0478.lateralMagnificationAt}
  \uses{def:physics:phyx-mini-0478:phyxminiproblems-problemphyxmini0478-lengthmagnitude, def:physics:phyx-mini-0478:phyxminiproblems-problemphyxmini0478-lengthincentimeters, def:physics:phyx-mini-0478:phyxminiproblems-problemphyxmini0478-mirrormagnificationexperiment}
  The graph's dimensionless magnification at a physical object distance
\end{definition}
\begin{proof}
  This is the declared model definition of lateral Magnification At.
\end{proof}
\begin{definition}[Matches Spherical Mirror Scenario]
  \label{def:physics:phyx-mini-0478:phyxminiproblems-problemphyxmini0478-matchessphericalmirrorscenario}
  \lean{PhyXMiniProblems.ProblemPhyXMini0478.MatchesSphericalMirrorScenario}
  \uses{def:physics:phyx-mini-0478:phyxminiproblems-problemphyxmini0478-mirrormagnificationexperiment}
  The spherical, on-axis, paraxial scenario described by the problem
\end{definition}
\begin{proof}
  This is the declared model definition of Matches Spherical Mirror Scenario.
\end{proof}
\begin{definition}[Matches Textually Described Magnification Graph]
  \label{def:physics:phyx-mini-0478:phyxminiproblems-problemphyxmini0478-matchestextuallydescribedmagnificationgraph}
  \lean{PhyXMiniProblems.ProblemPhyXMini0478.MatchesTextuallyDescribedMagnificationGraph}
  \uses{def:physics:phyx-mini-0478:phyxminiproblems-problemphyxmini0478-lengthincentimeters, def:physics:phyx-mini-0478:phyxminiproblems-problemphyxmini0478-mirrormagnificationexperiment}
  Axis labels, displayed ranges, and qualitative properties of the intended magnification graph. These fields do not state the requested value at 21 cm
\end{definition}
\begin{proof}
  This is the declared model definition of Matches Textually Described Magnification Graph.
\end{proof}
\begin{definition}[Has Intended Graph Calibration Readout]
  \label{def:physics:phyx-mini-0478:phyxminiproblems-problemphyxmini0478-hasintendedgraphcalibrationreadout}
  \lean{PhyXMiniProblems.ProblemPhyXMini0478.HasIntendedGraphCalibrationReadout}
  \uses{def:physics:phyx-mini-0478:phyxminiproblems-problemphyxmini0478-mirrormagnificationexperiment, def:physics:phyx-mini-0478:phyxminiproblems-problemphyxmini0478-lateralmagnificationat}
  The calibration point read from the intended curve: at the scale mark p\_s = 10 cm, the magnification is the marked value 0.5. This is a figure-derived input at 10 cm, not the requested conclusion at 21 cm
\end{definition}
\begin{proof}
  This is the declared model definition of Has Intended Graph Calibration Readout.
\end{proof}
\begin{definition}[Matches Question Readout]
  \label{def:physics:phyx-mini-0478:phyxminiproblems-problemphyxmini0478-matchesquestionreadout}
  \lean{PhyXMiniProblems.ProblemPhyXMini0478.MatchesQuestionReadout}
  \uses{def:physics:phyx-mini-0478:phyxminiproblems-problemphyxmini0478-lengthincentimeters, def:physics:phyx-mini-0478:phyxminiproblems-problemphyxmini0478-mirrormagnificationexperiment}
  The object-distance readout named by the question
\end{definition}
\begin{proof}
  This is the declared model definition of Matches Question Readout.
\end{proof}
\begin{definition}[Has Physical Mirror Parameters]
  \label{def:physics:phyx-mini-0478:phyxminiproblems-problemphyxmini0478-hasphysicalmirrorparameters}
  \lean{PhyXMiniProblems.ProblemPhyXMini0478.HasPhysicalMirrorParameters}
  \uses{def:physics:phyx-mini-0478:phyxminiproblems-problemphyxmini0478-lengthmagnitude, def:physics:phyx-mini-0478:phyxminiproblems-problemphyxmini0478-lengthmagnitudereadout, def:physics:phyx-mini-0478:phyxminiproblems-problemphyxmini0478-signedlengthreadout, def:physics:phyx-mini-0478:phyxminiproblems-problemphyxmini0478-mirrormagnificationexperiment}
  Positivity and nondegeneracy conditions for the physical branch. They assign no numerical value to f, q, or the requested magnification
\end{definition}
\begin{proof}
  This is the declared model definition of Has Physical Mirror Parameters.
\end{proof}
\begin{definition}[Satisfies Paraxial Spherical Mirror Laws]
  \label{def:physics:phyx-mini-0478:phyxminiproblems-problemphyxmini0478-satisfiesparaxialsphericalmirrorlaws}
  \lean{PhyXMiniProblems.ProblemPhyXMini0478.SatisfiesParaxialSphericalMirrorLaws}
  \uses{def:physics:phyx-mini-0478:phyxminiproblems-problemphyxmini0478-lengthmagnitude, def:physics:phyx-mini-0478:phyxminiproblems-problemphyxmini0478-lengthmagnitudereadout, def:physics:phyx-mini-0478:phyxminiproblems-problemphyxmini0478-signedlengthreadout, def:physics:phyx-mini-0478:phyxminiproblems-problemphyxmini0478-mirrormagnificationexperiment, def:physics:phyx-mini-0478:phyxminiproblems-problemphyxmini0478-lateralmagnificationat}
  The Gaussian spherical-mirror equation and the signed transverse- magnification law, imposed at every positive object distance and in every length unit: 1 / f = 1 / p + 1 / q; m = -q / p. These are governing laws and contain no numerical value at the queried distance
\end{definition}
\begin{proof}
  This is the declared model definition of Satisfies Paraxial Spherical Mirror Laws.
\end{proof}
\begin{lemma}[image Distance At Scale In Centimeters eq neg five]
  \label{lem:physics:phyx-mini-0478:phyxminiproblems-problemphyxmini0478-imagedistanceatscaleincentimeters-eq-neg-five}
  \lean{PhyXMiniProblems.ProblemPhyXMini0478.imageDistanceAtScaleInCentimeters_eq_neg_five}
  \uses{def:physics:phyx-mini-0478:phyxminiproblems-problemphyxmini0478-signedlengthincentimeters, def:physics:phyx-mini-0478:phyxminiproblems-problemphyxmini0478-mirrormagnificationexperiment, def:physics:phyx-mini-0478:phyxminiproblems-problemphyxmini0478-matchestextuallydescribedmagnificationgraph, def:physics:phyx-mini-0478:phyxminiproblems-problemphyxmini0478-hasintendedgraphcalibrationreadout, def:physics:phyx-mini-0478:phyxminiproblems-problemphyxmini0478-hasphysicalmirrorparameters, def:physics:phyx-mini-0478:phyxminiproblems-problemphyxmini0478-satisfiesparaxialsphericalmirrorlaws}
  The 10 cm, m = 0.5 calibration gives q = -5 cm
\end{lemma}
\begin{proof}
  The conclusion follows from the governing relations and figure readouts named in the statement of image Distance At Scale In Centimeters eq neg five.
\end{proof}
\begin{lemma}[focal Length In Centimeters eq neg ten]
  \label{lem:physics:phyx-mini-0478:phyxminiproblems-problemphyxmini0478-focallengthincentimeters-eq-neg-ten}
  \lean{PhyXMiniProblems.ProblemPhyXMini0478.focalLengthInCentimeters_eq_neg_ten}
  \uses{def:physics:phyx-mini-0478:phyxminiproblems-problemphyxmini0478-signedlengthincentimeters, def:physics:phyx-mini-0478:phyxminiproblems-problemphyxmini0478-mirrormagnificationexperiment, def:physics:phyx-mini-0478:phyxminiproblems-problemphyxmini0478-matchestextuallydescribedmagnificationgraph, def:physics:phyx-mini-0478:phyxminiproblems-problemphyxmini0478-hasintendedgraphcalibrationreadout, def:physics:phyx-mini-0478:phyxminiproblems-problemphyxmini0478-hasphysicalmirrorparameters, def:physics:phyx-mini-0478:phyxminiproblems-problemphyxmini0478-satisfiesparaxialsphericalmirrorlaws}
  The scale-point calibration and mirror equation give f = -10 cm
\end{lemma}
\begin{proof}
  The conclusion follows from the governing relations and figure readouts named in the statement of focal Length In Centimeters eq neg ten.
\end{proof}
\begin{lemma}[image Distance At Query In Centimeters eq neg two hundred ten div thirty one]
  \label{lem:physics:phyx-mini-0478:phyxminiproblems-problemphyxmini0478-imagedistanceatqueryincentimeters-eq-neg-two-hundred-ten-div-thirty-one}
  \lean{PhyXMiniProblems.ProblemPhyXMini0478.imageDistanceAtQueryInCentimeters_eq_neg_two_hundred_ten_div_thirty_one}
  \uses{def:physics:phyx-mini-0478:phyxminiproblems-problemphyxmini0478-signedlengthincentimeters, def:physics:phyx-mini-0478:phyxminiproblems-problemphyxmini0478-mirrormagnificationexperiment, def:physics:phyx-mini-0478:phyxminiproblems-problemphyxmini0478-matchestextuallydescribedmagnificationgraph, def:physics:phyx-mini-0478:phyxminiproblems-problemphyxmini0478-hasintendedgraphcalibrationreadout, def:physics:phyx-mini-0478:phyxminiproblems-problemphyxmini0478-matchesquestionreadout, def:physics:phyx-mini-0478:phyxminiproblems-problemphyxmini0478-hasphysicalmirrorparameters, def:physics:phyx-mini-0478:phyxminiproblems-problemphyxmini0478-satisfiesparaxialsphericalmirrorlaws}
  At p = 21 cm, the mirror equation gives q = -210/31 cm
\end{lemma}
\begin{proof}
  The conclusion follows from the governing relations and figure readouts named in the statement of image Distance At Query In Centimeters eq neg two hundred ten div thirty one.
\end{proof}
\begin{definition}[Answer Choice]
  \label{def:physics:phyx-mini-0478:phyxminiproblems-problemphyxmini0478-answerchoice}
  \lean{PhyXMiniProblems.ProblemPhyXMini0478.AnswerChoice}
  Labels of the four answer choices printed in the source problem
\end{definition}
\begin{proof}
  This is the declared model definition of Answer Choice.
\end{proof}
\begin{definition}[displayed Magnification]
  \label{def:physics:phyx-mini-0478:phyxminiproblems-problemphyxmini0478-displayedmagnification}
  \lean{PhyXMiniProblems.ProblemPhyXMini0478.displayedMagnification}
  \uses{def:physics:phyx-mini-0478:phyxminiproblems-problemphyxmini0478-answerchoice}
  Dimensionless magnification printed beside each answer choice
\end{definition}
\begin{proof}
  This is the declared model definition of displayed Magnification.
\end{proof}
\begin{definition}[recorded Dataset Answer]
  \label{def:physics:phyx-mini-0478:phyxminiproblems-problemphyxmini0478-recordeddatasetanswer}
  \lean{PhyXMiniProblems.ProblemPhyXMini0478.recordedDatasetAnswer}
  \uses{def:physics:phyx-mini-0478:phyxminiproblems-problemphyxmini0478-answerchoice}
  The answer label recorded by the source dataset, retained as metadata
\end{definition}
\begin{proof}
  This is the declared model definition of recorded Dataset Answer.
\end{proof}
\begin{definition}[Matches Displayed Magnification]
  \label{def:physics:phyx-mini-0478:phyxminiproblems-problemphyxmini0478-matchesdisplayedmagnification}
  \lean{PhyXMiniProblems.ProblemPhyXMini0478.MatchesDisplayedMagnification}
  \uses{def:physics:phyx-mini-0478:phyxminiproblems-problemphyxmini0478-answerchoice, def:physics:phyx-mini-0478:phyxminiproblems-problemphyxmini0478-displayedmagnification}
  Agreement with a displayed magnification to the nearest hundredth
\end{definition}
\begin{proof}
  This is the declared model definition of Matches Displayed Magnification.
\end{proof}
\begin{definition}[Source scene]
  \label{def:physics:phyx-mini-0478:phyxminiproblems-problemphyxmini0478-sourcescene}
  \lean{PhyXMiniProblems.ProblemPhyXMini0478.SourceScene}
  The source package contains either the prose's spherical-mirror graph or the supplied raster's unrelated electric-field vector diagram.
\end{definition}
\begin{proof}
  This is the declared source-scene classification.
\end{proof}
\begin{definition}[Evidence policy]
  \label{def:physics:phyx-mini-0478:phyxminiproblems-problemphyxmini0478-evidencepolicy}
  \lean{PhyXMiniProblems.ProblemPhyXMini0478.EvidencePolicy}
  Because the raster conflicts with the prose, optical setup information is taken only from the prose and the raster is retained solely as an audited mismatch.
\end{definition}
\begin{proof}
  This is the declared evidence policy.
\end{proof}
\begin{definition}[Raster point]
  \label{def:physics:phyx-mini-0478:phyxminiproblems-problemphyxmini0478-rasterpoint}
  \lean{PhyXMiniProblems.ProblemPhyXMini0478.RasterPoint}
  These are the six numbered points visible in the supplied vector-diagram raster.
\end{definition}
\begin{proof}
  This is the finite point-label classification read from the raster.
\end{proof}
\begin{definition}[Supplied raster audit]
  \label{def:physics:phyx-mini-0478:phyxminiproblems-problemphyxmini0478-suppliedrasteraudit}
  \lean{PhyXMiniProblems.ProblemPhyXMini0478.SuppliedRasterAudit}
  \uses{def:physics:phyx-mini-0478:phyxminiproblems-problemphyxmini0478-sourcescene, def:physics:phyx-mini-0478:phyxminiproblems-problemphyxmini0478-rasterpoint}
  The audit records the raster's scene and the optional integer component label attached to each numbered point.
\end{definition}
\begin{proof}
  This is the data structure for the literal raster transcription.
\end{proof}
\begin{definition}[Records supplied raster mismatch]
  \label{def:physics:phyx-mini-0478:phyxminiproblems-problemphyxmini0478-recordssuppliedrastermismatch}
  \lean{PhyXMiniProblems.ProblemPhyXMini0478.RecordsSuppliedRasterMismatch}
  \uses{def:physics:phyx-mini-0478:phyxminiproblems-problemphyxmini0478-suppliedrasteraudit, def:physics:phyx-mini-0478:phyxminiproblems-problemphyxmini0478-sourcescene, def:physics:phyx-mini-0478:phyxminiproblems-problemphyxmini0478-rasterpoint}
  The audit identifies the raster as the electric-field scene and transcribes the five printed component pairs, with no arrow label at point six.
\end{definition}
\begin{proof}
  This is the conjunction of the six literal raster observations.
\end{proof}
\begin{definition}[Transverse imaging remainder]
  \label{def:physics:phyx-mini-0478:phyxminiproblems-problemphyxmini0478-transverseimagingremainderincentimeters}
  \lean{PhyXMiniProblems.ProblemPhyXMini0478.transverseImagingRemainderInCentimeters}
  \uses{def:physics:phyx-mini-0478:phyxminiproblems-problemphyxmini0478-mirrormagnificationexperiment, def:physics:phyx-mini-0478:phyxminiproblems-problemphyxmini0478-signedlengthincentimeters, def:physics:phyx-mini-0478:phyxminiproblems-problemphyxmini0478-lengthincentimeters}
  At fixed object distance, this is the transverse image-height response minus its first-order signed magnification term.
\end{definition}
\begin{proof}
  This is the defining remainder for the local transverse imaging law.
\end{proof}
\begin{lemma}[Gaussian mirror relation at the query]
  \label{lem:physics:phyx-mini-0478:phyxminiproblems-problemphyxmini0478-gaussianmirrorrelationatquery}
  \lean{PhyXMiniProblems.ProblemPhyXMini0478.gaussianMirrorRelationAtQuery}
  \uses{def:physics:phyx-mini-0478:phyxminiproblems-problemphyxmini0478-mirrormagnificationexperiment, def:physics:phyx-mini-0478:phyxminiproblems-problemphyxmini0478-matchesquestionreadout, def:physics:phyx-mini-0478:phyxminiproblems-problemphyxmini0478-hasphysicalmirrorparameters, def:physics:phyx-mini-0478:phyxminiproblems-problemphyxmini0478-satisfiesparaxialsphericalmirrorlaws, def:physics:phyx-mini-0478:phyxminiproblems-problemphyxmini0478-signedlengthincentimeters, def:physics:phyx-mini-0478:phyxminiproblems-problemphyxmini0478-lengthincentimeters}
  At the queried object distance, the focal length times the sum of object and signed image distance equals their product.
\end{lemma}
\begin{proof}
  Differentiate the exact reflection residual at the optical axis, use its local vanishing law, and clear the nonzero physical denominators.  This yields the division-free Gaussian relation.
\end{proof}
\begin{lemma}[Transverse magnification relation at the query]
  \label{lem:physics:phyx-mini-0478:phyxminiproblems-problemphyxmini0478-transversemagnificationrelationatquery}
  \lean{PhyXMiniProblems.ProblemPhyXMini0478.transverseMagnificationRelationAtQuery}
  \uses{def:physics:phyx-mini-0478:phyxminiproblems-problemphyxmini0478-mirrormagnificationexperiment, def:physics:phyx-mini-0478:phyxminiproblems-problemphyxmini0478-lateralmagnificationat, def:physics:phyx-mini-0478:phyxminiproblems-problemphyxmini0478-matchesquestionreadout, def:physics:phyx-mini-0478:phyxminiproblems-problemphyxmini0478-hasphysicalmirrorparameters, def:physics:phyx-mini-0478:phyxminiproblems-problemphyxmini0478-satisfiesparaxialsphericalmirrorlaws, def:physics:phyx-mini-0478:phyxminiproblems-problemphyxmini0478-transverseimagingremainderincentimeters, def:physics:phyx-mini-0478:phyxminiproblems-problemphyxmini0478-signedlengthincentimeters, def:physics:phyx-mini-0478:phyxminiproblems-problemphyxmini0478-lengthincentimeters}
  At the queried distance, lateral magnification times object distance equals minus the signed image distance.
\end{lemma}
\begin{proof}
  Differentiate the transverse image-height decomposition at zero object height.  The remainder derivative is zero and the image-height derivative is the defined graph magnification, so clearing the positive object distance gives the result.
\end{proof}
% --- Archon declaration topology end ---
% --- Archon physics formalization source end ---
